% answers_short.tex —— 2009 年·普通高等学校招生考试·江西卷·文科 · 简短答案速查
% 由 tools/make_short_answers.py 自 answers.tex 的 \gksolhead/\gkanswer 提取，
% 并与 2009/2009_数学真题_江西文_答案要点.md 逐题比对后生成，勿手改（改源 MD 后重跑）。

\begin{gknotice}
\textbf{使用说明}：本节只给最终答案便于快速核对，完整推导见解析版。答案为编写者独立推导，\textbf{非官方标准答案}。
\end{gknotice}

\gkpart{选择题}
\begin{center}\begin{tabular}{ccccccccccccc}
\toprule
题号 & 1 & 2 & 3 & 4 & 5 & 6 & 7 & 8 & 9 & 10 & 11 & 12 \\
\midrule
答案 & A & D & B & A & C & C & B & C & C & D & B & A \\
\bottomrule
\end{tabular}\end{center}

\gkpart{填空题}
\begin{itemize}[leftmargin=2.6em]
  \item[\textbf{13.}] $0$。
  \item[\textbf{14.}] $\dfrac{8\sqrt{6\pi}}{\pi}$。
  \item[\textbf{15.}] $\dfrac{\sqrt3}{2}$。
  \item[\textbf{16.}] A、B、C。
\end{itemize}

\gkpart{解答题}
\begin{itemize}[leftmargin=2.6em]
  \item[\textbf{17.}] （1）$-\dfrac34$；（2）$a<2$ 或 $a>\dfrac52$。
  \item[\textbf{18.}] （1）$\dfrac1{64}$；（2）$\dfrac{11}{32}$。
  \item[\textbf{19.}] （1）$C=\dfrac{\pi}{4}$；（2）$a=\sqrt2$，$b=\sqrt3+1$，$c=2$。
  \item[\textbf{20.}] （1）证明见解析；（2）$\arccos\dfrac13$；（3）$\sqrt2$。
  \item[\textbf{21.}] （1）$n=3k$ 时 $S_n=\dfrac{k(9k+4)}{2}$，$n=3k+1$ 时 $S_n=-\dfrac{2k+1}{2}$，$n=3k+2$ 时 $S_n=-\dfrac{9k^2+14k+5}{2}$（$k$ 为非负整数，$n=3k$ 时 $k\geqslant1$）；（2）$T_n=\dfrac83-\dfrac{9n+16}{6\cdot4^n}$。
  \item[\textbf{22.}] （1）$r=\dfrac23$；（2）证明见解析。
\end{itemize}
